\documentclass[11pt]{article}
\usepackage{amsmath,amssymb,amsthm,hyperref}
\newtheorem{theorem}{Theorem}\newtheorem{proposition}[theorem]{Proposition}
\title{Scheduled event ideals, native cuts, and a scheduler-invariant size clock}
\date{16 November 2026}
\begin{document}\maketitle
\paragraph{Audit and target.} The written extension passed independent adversarial review and exact finite diagnostics at its stated forest interface. No novelty, compilation, or formal-verification claim is made. This replaces the all-frontier round interface by scheduled finite enabled antichains, retaining pending events and chosen conditional laws. Tensor coefficients are retained. The main assembly theorem covers arbitrary stopping ideals between atomic draws, a full marked ideal grade, and a causally constructed node-count clock. Independent terminal evaluation is scheduler-invariant even for unfair nonstalling schedulers, whereas recurrent prefixes and shared-control evaluation are not. The distinction is proved, not hidden behind a scheduler-invariance slogan.

\section{The event interface and full data}
Use the reviewed tensor-causal source's primitive $f:A+U\to F_C(B+U)$, countable raw rows, augmented context coefficients $E_{x,w}\to k$, and independently weighted full content collection $C$. Fix a finite number $m$ of ordered external roots, with inputs in $A^m$. Each event is one original atomic primitive call. Its raw result determines its ordered children: settled $B$ leaves never execute, and continuing $U$ children become enabled primitive calls immediately after their parent returns. In particular a child can execute while a sibling of its parent remains pending; there is no all-sibling barrier. Event identifiers are stable planar root/node addresses. The full terminal grade $R_U^{(m)}C$ is the entire ordered $m$-forest of finite nonempty $C$-trees, with internal $U$ labels and all actual finite content tensor factors, including every unused forest. For $m=0$ it is the singleton empty forest with unit coefficient and no positions.

A raw past records every executed node and its raw outcome, settled leaves, all pending typed inputs, and the ordered sequence of scopes already sampled. Its executed node set is an ancestor-closed ideal in the resulting partial forest. The enabled set is the current finite frontier. A scheduler $s$, defined on every syntactic nonterminal raw past including zero-probability ones, chooses a nonempty subset $S(h)$ of that frontier. These selected events form an antichain. All others remain pending. The scheduler is deterministic given the full retained past. A complete control supplies a normalized joint law
\[
 q_h\text{ on }\prod_{v\in S(h)}W_{x_v(h)},\qquad
 (\pi_v)_*q_h=p_{x_v(h)}                                   \tag{1}
\]
conditional on the entire past. This is an atomic selected-scope draw. A legal cut occurs between such draws, never inside one. The selected scopes may have shared innovation even though their input labels are causally known before the draw. Their tensor coefficients are those of the original chosen-control open crossing; no product-coefficient universal property is presumed.

The evaluated raw histories are precisely the syntactic histories consistent with the chosen scheduler decisions, including every zero-probability draw at those decisions. Incompatible alternative plans are not inserted as artificial zero raw histories. The full grade nevertheless retains all legal plans, and the complete scheduler/control table is required to be defined on all legal raw pasts, including pasts incompatible with its earlier decisions; those extra contingent data remain part of the control index.

A terminal history has no pending event. Terminal histories are absorbed as themselves at later bookkeeping horizons, not padded by artificial raw idle events. Every finite history has finitely many executed nodes and finitely many pending calls. Countably many finite histories are reachable from each external input, because each selected scope is finite and all primitive raw rows are countable. Actual runs may be infinite; no scalar measure on infinite runs or accumulated infinite coefficient is constructed here.

\section{Scheduled finite prefixes and full plan grades}
A prefix object after $n$ selected-scope steps retains every syntactically possible active history of exact step $n$ and every already terminal history of at most $n$ steps. Its weight is the product of its conditional scope probabilities and its context coefficient is the tensor over its executed nodes. It has mass one by induction: each active history distributes its mass by (1), and a terminal history persists exactly once.

Full grades must retain enough information to make truncation defined independently of raw realizability. Let $Q_n$ have shapes consisting of a finite partial ordered $C$-forest $p$, all internal and active-edge $U$ labels, and a legal ordered scope plan $\pi=(S_1,\ldots,S_k)$ of its executed nodes. The scopes are nonempty, execute only enabled nodes, and partition the executed set. If the prefix is active, require $k=n$; if it is terminal, require $k\leq n$. At $n=0$ the shape is the ordered $m$-seed forest, with its external root inputs retained in raw state rather than placed in the grade. For $m=0$ this seed is already terminal. All legal plans are included, including plans never chosen by the actual scheduler and all unused shapes. Variable positions are only settled $B$ leaves; active holes and their labels are constants. Content coefficients are the finite node tensors.

For prefix truncation assume coherent content augmentations and the enlarged backward-position grade maps of the reviewed prefix theorem. There is a full map $Q_{n+1}\to Q_n$: an already terminal plan with $k\leq n$ is unchanged; otherwise remove its last scope, delete those executed nodes, and expose their incoming $U$ edges as pending holes. If the removed scope contains a root, recover its unexecuted external root hole, not a $U$-label. Its root input remains in raw state. The backward position map includes retained settled leaves, and the content arrow augments deleted node factors. A last-scope node has no executed descendant, since its descendants were not enabled before that draw. Therefore this operation is defined on all full plan shapes. Retaining a plan is material: a bare non-level-complete partial tree need not determine which of several enabled antichains was executed last.

\begin{theorem}[Scheduled native prefix tower]
The normalized prefix objects at full grades $Q_n$ and raw last-scope truncation form a native inverse sequence. Context arrows delete new factors by augmentation. Full shape maps, backward positions and content arrows are defined on all unused plan shapes. Given a primitive native comparison inducing a history map $H$, assume the local scheduler and control equations
\[
 S'(Hh)=H(S(h)),\qquad (r_{S(h)})_*q_h=q'_{Hh}.             \tag{2}
\]
Then the induced prefix maps are actual native maps commuting with truncation. Tensor-to-product comparisons at each retained executed node set form a natural map to the analogous product prefix tower.
\end{theorem}
\begin{proof}
The fibre of raw truncation above an active prefix is precisely its selected-scope draw set; its total successor mass is the prefix mass times $\sum q_h=1$. Above an already terminal prefix the fibre is a singleton. This proves exact pushforward, including zero prefixes. The truncated body is exactly the full plan truncation just defined. Retained settled positions are included backward, and deleting context/content factors uses actual augmentations. Multiple deletions commute by symmetric monoidal coherence and monoidality of augmentations. Thus these are native arrows, not effect-only equalities.

The first equation in (2) ensures corresponding histories execute the same addressed events at the next step; primitive body squares identify their new ordered frontiers. The second equation gives the inner draw-fibre sum, and induction on steps gives exact prefix pushforward. Actual primitive coefficient arrows tensor over corresponding executed nodes; augmentation preservation makes them commute with truncation. The full collection and state maps send all plan shapes addresswise, preserve their plans, and commute with their full grade truncations. Finally the tensor-to-product comparison commutes with any finite node deletion and with all node maps by its proved coefficient calculus. Raw maps and full grade maps for that comparison are identities, proving the tower square.
\end{proof}

For $m=1$, the pinned older whole-frontier prefix tower embeds; the same statement holds for its explicitly ordered finite-forest analogue. The embedding is given by assigning each of its full level-complete shapes the canonical plan that executes every current frontier. This is a full-grade inclusion into $Q_n$, with unchanged raw histories and coefficient arrows when the chosen scheduler is all-frontier. It is not an isomorphism onto the larger unused-complete $Q_n$. At termination the common unmarked grade $R_U^{(m)}C$ will remove this discrepancy through the explicitly declared forget-plan map, not by a hidden support restriction.

\section{Terminal evaluation and stopping event ideals}
The tensor terminating process has every finite terminal scheduled raw history $t$ as a witness, weight the product of its scope probabilities, and coefficient the tensor of its executed node coefficients. Its body forgets the schedule but retains the full terminal forest and ordered $B$ leaves at grade $R_U^{(m)}C$. Terminal mass is at most one, by finite-horizon conservation. Zero-weight terminal histories remain. The forget-plan body is not an assertion that all schedules or shared laws have equal weights.

Fix a deterministic stopping rule on ALL legal raw pasts, including pasts incompatible with the chosen scheduler's earlier decisions, bounded by $N$ selected-scope steps and always stopping earlier at termination. Its decisions occur only at scope boundaries. A cut retains the entire raw past, its executed ideal $J$, settled leaves, pending frontier and complete contingent suffix policy. The suffix operates jointly over that retained frontier, not as independent branch continuations. Its typed boundary can mix unexecuted external $A$-roots with internal $U$-children; the actual input and its tag remain at each address, rather than recasting all pending calls as $U$-holes. For a terminal forest $t$, the full marked cut grade has every pair $(t,J)$ with $J$ an ancestor-closed node subset, including the empty ideal, the entire node set, and all unused intermediate ideals. Its coefficient is the same full node tensor. This grade does not assert that every unused ideal can be reached by the chosen policy: actual cuts must occur between its atomic scopes. Forgetting the ideal is a full generally noninvertible grade map to $R_U^{(m)}C$. Multiple cuts retain chains $J_1\subseteq\cdots\subseteq J_r$.

\begin{theorem}[Stopping-ideal native assembly]
A complete scheduled control, with tables on all legal pasts, decomposes uniquely into its prefix control and one complete suffix control at EVERY legal stopping past, including zero pasts and scheduler-incompatible pasts. Every terminal history decomposes uniquely into that past and a remaining scheduled suffix forest. Its weight factors as prefix weight times conditional suffix weight, and its actual tensor coefficient factors over the two disjoint node sets. This gives an invertible context assembly paired with the full directed marked-grade forgetting map. Repeated cuts and the native comparisons (2) commute with assembly when the stopping decisions commute with $H$. The tensor-to-product comparison also commutes with this assembly.
\end{theorem}
\begin{proof}
Every nonterminal raw history lies strictly before its unique first stopping past or extends precisely one such past. The scheduler decision and conditional table at that history therefore belong to exactly one part of the decomposition, even when it is off-policy or its prefix mass is zero. Prefix and suffix table domains cover all legal histories, whereas their evaluated raw laws use only scheduler-compatible histories. Conversely those complete tables reconstruct the whole control. No conditional distribution is inferred by division at zero histories.

A terminal history's finite scope sequence has a unique stopping split. Its prefix and suffix retain the full raw past, so concatenation and grafting are inverse on raw histories. Splitting the product of scope probabilities proves the weight equation. The executed nodes split into $J$ and its complement; actual tensor associators and symmetries identify the total coefficient with the tensor of the two parts. The full content coefficients similarly factor, and the marked grade remembers $J$ on every unused shape before forgetting it. For repeated cuts, every route concatenates the same raw sequence and partitions the same addressed nodes, giving the same actual maps by coherence. Compatible native comparisons preserve the stopping split and have the same node maps. The tensor-to-product comparison commutes with arbitrary finite partitions, so its assembly square commutes too.
\end{proof}

The reconstruction statement does not apply to an arbitrary marked ideal that splits a sampled scope or fails to be an initial block prefix of the raw plan. Those unused full-grade points do not create admitted stopping policies.

This admits genuinely non-level cuts. A root can enable three events; a scheduler jointly draws a shared fair bit for two of them while leaving the third pending, then stops. Its executed ideal contains the root and precisely those two children. The cut is legal and retains their correlation, but it is not a whole-frontier depth cut. Cutting between the two members of that already chosen atomic draw remains illegal under (1).

\section{Independent rescheduling and its precise naturality}
A complete syntactic terminal raw forest records every primitive raw witness, input and ordered node address, but no scope plan. Fix a deterministic nonstalling scheduler as above.

\begin{theorem}[Finite terminal-tree reconstruction]
Every such finite terminal raw forest has exactly one scheduled execution under the scheduler. Its number of selected-scope steps is at most its number of executed nodes. Consequently, under independent controls $q_h=\prod_{v\in S(h)}p_{x_v(h)}$, all deterministic nonstalling schedulers give natively isomorphic terminating processes at identity full grade $R_U^{(m)}C$. Their law at a raw forest is the product of all primitive node probabilities and their context coefficient is its original node tensor. This includes equality of total finite-termination masses. No fairness assumption is needed for this finite-terminal claim.
\end{theorem}
\begin{proof}
Given the forest, simulate the scheduler from its seed past. Every enabled event is an as-yet-unexecuted node of that forest. The forest supplies the raw witness at every event selected by the scheduler, so the next raw past is determined uniquely. Every selected scope is nonempty and executes new nodes, reducing the finite number of unexecuted nodes. No enabled event can lie outside the given complete forest, since all continuing children and their raw subtrees are already specified there. Simulation must therefore terminate within the number of nodes, and has no choices. This includes all zero raw forests. Conversely forgetting a scheduled terminal history gives that forest and reconstructs the same history.

For independent controls the finite product over scopes is exactly the product over their disjoint node sets. It is independent of the plan. The coefficient tensor also has those same factors and uses its actual addressed symmetry/associativity map. The terminal body and complete global grade are unchanged. Thus the raw bijection with these actual coefficient isomorphisms gives native inverse arrows, not only an equality of probabilities. Summing equal weights over all finite terminal forests gives the same termination mass.
\end{proof}

Moreover, between different primitive presentations with independent controls, \emph{every} original arity-preserving primitive native comparison induces exact native terminal transport, even if the two deterministic schedulers do not satisfy the scheduler equation (2). Work on complete raw forests: the primitive body squares give an address-preserving forest map, and tensor the actual node maps. Exact pushforward follows by induction on a finite target forest. Sum over source root witnesses above its root witness; for each, sum independently over its source subtrees. The inductive subtree sums are the same target subtree weights for every source child input mapping to the given target input. Their product can therefore be factored from the root sum, whose value is the target root weight by primitive exact pushforward. Apply the same reasoning to the finite ordered roots. Countable nonnegative sums justify the fibre regrouping. The full content map is $R_u^{(m)}\alpha$ on every shape. This naturality occurs after terminal schedule forgetting; without (2) it does not assert a map between the two scheduled prefix towers.

\section{A causally generated node-count clock}
Step-count first-hit layers need not agree across schedulers. A batch of two enabled events can execute in one selected scope or two, even with independent laws. The traced terminating outputs agree, but there need not be any unclosed fixed-step-clock comparison.

For this clock theorem assume $A,B,U$ countable as in the pinned countable-endpoint trace ambient; the prefix and terminating constructions above require only small endpoints with countable reachable raw histories per external input. Use instead the node-count clock $n\in\mathbb N$. At row $(a,n)$, perform the finite scheduled prefix construction for at most $n$ selected-scope steps, absorbing terminal histories. Apply the independently supplied typed restriction to those terminal histories whose total executed node count is exactly $n$. Retype its full terminal-plan subcollection and apply the full forget-plan map into the entire $R_U^{(m)}C$. Return clock output $n$. This construction requires no already evaluated infinite terminating law. For $m>0$, the zero-clock row is empty. For $m=0$, it is the single empty forest of weight one; all positive-clock rows are empty. All selected terminal zero witnesses remain present.

\begin{theorem}[Size-clock assembly and rescheduling]
The node-count clock is admitted by the countable subprobability trace and traces to the scheduled terminating evaluator. It commutes with native comparisons (2) and the tensor-to-product comparison. On independent controls, even arbitrary deterministic nonstalling scheduler changes and all original primitive native comparisons induce native maps between its rows at the \emph{same node-count clock}. Its trace agrees natively with the earlier first-step-hit clock result, without requiring an unclosed map from step count to node count.
\end{theorem}
\begin{proof}
A terminal history with $n$ nodes uses at most $n$ nonempty scopes, so it appears in the $n$-step prefix construction. Absorption records it once. Restriction therefore selects exactly all terminal histories of node count $n$, not just those terminating at step $n$. Its mass is at most one. Every finite terminal history has a unique finite node count, so the sum of these row masses is its total terminating mass, at most one. Countable trace is therefore admitted. Its raw bijection to all terminal histories preserves every coefficient, body and full grade. The first-step-hit clock gives another disjoint indexing of the same histories, hence the same traced result. A step count alone does not determine a node count; no unsupported unclosed clock reindexing is used.

Comparisons (2) preserve addressed nodes and transport the finite prefixes and selection predicate, giving fixed-clock native maps. Tensor-to-product comparison changes no raw count or full grade, and commutes with restriction and trace exactly as before. For arbitrary independent rescheduling and primitive comparisons, the complete-forest proof in the preceding section is node-count preserving. Every source forest over a target forest has the same addressed node count. Its exact fibre sum thus holds within each clock row separately. The context arrows tensor the same node maps, and full grade maps are $R_u^{(m)}\alpha$. This proves fixed-size-clock native transport directly, not by pretending the incompatible scheduled-prefix maps exist.
\end{proof}

\section{Failure boundaries and assembly interface}
Shared-control scheduler invariance is false. At two fair-bit terminal siblings, a simultaneous scope admits diagonal coupling. If the first outcome is revealed in a separate scope, the unchanged primitive fair marginal at the second scope forbids retaining that diagonal joint law. This is the established scope-splitting obstruction, now at an enabled-antichain interface. Schedule changes can be compared for shared controls only with separately supplied compatible conditional-law data; primitive marginal agreement alone does not supply it.

Independent terminal invariance does not imply recurrent prefix or liveness invariance. A root can spawn one perpetually continuing branch and one immediately terminating branch. A scheduler that always selects the continuing branch can leave the other event pending forever; a fair scheduler can execute the terminating branch first. Both have no finite complete terminal forest, but their settled-leaf prefix observations differ. Nonstalling is not fairness. No comparison of infinite execution laws, or theorem that fairness makes full prefix towers isomorphic, is claimed. The retained prefix tower and pending ports expose rather than erase this difference.

The source interface here has finitely many initial roots, finite primitive arities and countable raw rows. Its executions may be infinite while all finite prefixes have finite executed and pending sets. A countably infinite initial forest is a different interface requiring a grade for infinitely many pending constant holes; it is not silently included and cannot finish in finitely many finite scopes if every root must execute. No infinite tensor is supplied.

For finite CPS/open-port composition, the consumable mechanism is explicit: events are atomic original calls; child tasks become enabled immediately after the parent shape returns; a selected enabled antichain has a chosen conditional joint; cuts retain paid past and pending typed ports. The coefficient, native transport, event-ideal cut and size-clock arguments apply when its independently constructed full grade/body interface satisfies these address and truncation rules. The finite CPS grade remains its original unused-complete composite collection, such as $C\circ D$; this recurrence theorem does not replace it by $R_UC$ or assert those distinct full grades are isomorphic. Proving that finite syntax's grade/body realization is its author's separate obligation. The present construction proves the recurrent forest instance and states precisely what that realization must preserve.
\end{document}
